\section{Shell and Shell Commands}
\label{sec:shell}

Apollo-RTOS includes a Unix-like shell that executes commands sent over the serial bus.
The shell runs as a high-priority process that sleeps while waiting for input.

\begin{itemize}
    \item By default, all UART inputs are forwarded to the listener registered by
        \lstinline[style=myagc]|shell|. This listener buffers characters until a line
        break (\lstinline[style=myagc]|'\r'| or \lstinline[style=myagc]|'\n'|) is
        received, then wakes the shell process.
    \item Upon waking, \lstinline[style=myagc]|shell| parses the input to extract the
        command name, its arguments, background execution flag, and any output
        redirection to a file.
    \item It then searches the \lstinline[style=myagc]|.apps| section for a matching
    process, packages the arguments into an \lstinline[style=myagc]|args_t| structure
    (\autoref{ds:args}), and schedules the new process with this structure as its
        parameter.
    \item Finally, it clears the input buffer and returns to sleep.
\end{itemize}

Shell commands follow the syntax:
\begin{lstlisting}[style=mytty]
cmd [arguments] [> fn] [&]
\end{lstlisting}
For example, \lstinline[style=mytty]|echo hello, world > 10| writes
\lstinline[style=mytty]|hello, world| to the file \lstinline[style=mytty]|10|.

Available commands are listed using the \lstinline[style=myagc]|help| command:
\begin{lstlisting}[style=mytty]
>> help
pkill   htop    help    prof    accel   show
calc    clear   echo    cat     ls      touch
rm      heart   
\end{lstlisting}

Alongside familiar Unix-like commands, Apollo-RTOS provides a few custom commands to
demonstrate system features:

\paragraph*{\lstinline[style=myagc]|htop|} Displays running processes, waitlist tasks,
and open files:
\vspace{1em}
\begin{lstlisting}[style=mytty]
>> htop
boot level: flash
curr_proc: htop
  |  0. idle : (idle), [runnable]
  |  1. shell : (high), [asleep]
  |  2. display : (low), [runnable]
  |  3. accel_bg : (high), [runnable]
  |  4. htop : (medium), [running]
waitlist: NONE
\end{lstlisting}

\paragraph*{\lstinline[style=myagc]|accel|} Launches a simple accelerometer-based
compass on the LED matrix, pointing an arrow toward the ground. It can run in the
background or foreground:
\vspace{1em}
\begin{lstlisting}[style=mytty]
>> accel &
started [4] output to stdout

>> pkill accel
killed [4]

>> accel
x: 1, y: -5, z: 64
\end{lstlisting}

\vspace{1em}

\begin{figure}[H]
    \centering
    \begin{subfigure}[t]{0.3\textwidth}
        \centering
        \includegraphics[height=2in]{accel_bw1}
    \end{subfigure}%
    % 
    \begin{subfigure}[t]{0.3\textwidth}
        \centering
        \includegraphics[height=2in]{accel_bw3}
    \end{subfigure}%
    %
    \begin{subfigure}[t]{0.3\textwidth}
        \centering
        \includegraphics[height=2in]{accel_bw2}
    \end{subfigure}

    \caption{Result of \lstinline[style=myagc]|accel|}
\end{figure}


\newpage
\lstinputlisting[
    style=mytty,
    frame=single,
    basicstyle=\linespread{1}\ttfamily\codefontsize,
    caption={Sample workflow on the shell on \texttt{minicom}}
]{code/ex.sh}
